\begin{tikzpicture}[scale=1.0]
  % axes
  \draw[-Latex] (0,0) -- (8.4,0) node[right] {$c_1$};
  \draw[-Latex] (0,0) -- (0,7.0) node[above] {$c_2$};
  % feasible line c1+c2=6
  \draw[line width=1.2pt] (0,6) -- (6,0);
  \node[left] at (0,6) {$y$};
  \node[below] at (6,0) {$y$};
  % innermost indifference curve c1*c2=5 : passes through C=(1,5) and B=(5,1),
  % crossing the budget line on both sides.
  \draw[semithick,domain=0.85:5.85,smooth,variable=\x,samples=120]
        plot ({\x},{5/\x});
  % middle curve c1*c2=9 : tangent to the line at E=(3,3) (golden rule)
  \draw[semithick,domain=1.5:6.0,smooth,variable=\x,samples=120]
        plot ({\x},{9/\x});
  % outer curve c1*c2=14 : higher utility, infeasible, passes through D=(2.8,5)
  \draw[semithick,domain=2.0:7.6,smooth,variable=\x,samples=120]
        plot ({\x},{14/\x});
  % golden-rule allocation E
  \fill (3,3) circle (1.6pt);
  \node[below left=-1pt and 0pt] at (3,3) {$E$};
  \draw[densely dashed] (3,3) -- (3,0) node[below] {$c_1^*$};
  \draw[densely dashed] (3,3) -- (0,3) node[left] {$c_2^*$};
  % points C and B (the two crossings of the innermost curve with the line)
  \fill (1,5) circle (1.6pt);
  \node[above right=-1pt and 1pt] at (1,5) {$C$};
  \fill (5,1) circle (1.6pt);
  \node[above right=-1pt and 1pt] at (5,1) {$B$};
  % point D on the outer (unattainable) curve
  \fill (2.8,5) circle (1.6pt);
  \node[above right=-1pt and 1pt] at (2.8,5) {$D$};
  % annotation
  \node at (6.15,4.55) {The Golden Rule};
  \node at (6.6,3.95) {Allocation};
  \draw[-Latex] (5.2,4.2) -- (3.28,3.18);
\end{tikzpicture}
